\section{Methods}

\subsection{Prospective freeze}
Each result-bearing lane binds its subject, admissible information, candidates/baselines, endpoints, gates, tie-breaks, and terminal vocabulary before outcome inspection. Later corrections are logged rather than retroactively changing a gate. This makes a refutation or donor tie an admissible result rather than an incentive to move the goalposts.

\subsection{Donor first right of refusal}
A candidate receives scientific residual credit only after the strongest faithfully composable existing mechanism has been run or otherwise discharged under the same admitted information. A tie or donor win closes the candidate claim. Donor absorption is therefore a positive methodological terminal even though it yields zero candidate novelty credit.

\subsection{Typed failure and authority}
The workflow separates scientific state from orchestration state. Missing capabilities, failed tools, and malformed receipts are control conditions, not evidence against a hypothesis. Scientific/candidate receipts cannot by themselves grant novelty, adoption, promotion, merge, or global-stop authority. Reopening is tied to typed changes in assumptions, interfaces, representations, or evidence rather than to generic dissatisfaction with a negative result.

\subsection{Exactness and hostile controls}
Where a claim is finite and exact, the programme uses exhaustive enumeration, proof-carrying dynamic programming, independent witness reconstruction, or complete local machine checks. Protocols predeclare hostile controls designed to make an attractive shortcut fail. If the control does not bite where required, the world/lane is invalid rather than counted as a positive.

\subsection{Receipt and replay contract}
Result-bearing computations serialize canonical inputs, outputs, identifiers, and authority boundaries. Content-addressed or hash-bound receipts permit deterministic replay or independent recomputation. The contract establishes attribution and reproducibility; it does not establish the truth of a scientific inference without the corresponding design/referee.

\subsection{Case-study selection}
We analyze the complete R6 saturation/recovery arc and the N1--N4 closure families because they were defined by the programme before this paper synthesis and include heterogeneous terminals. The paper does not select only successful recoveries. Companion papers own the domain-specific claims.